\section{Introduction}
\label{sec:intro}

Large language model (LLM) agents increasingly depend on memory that outlives a single context
window---the facts they extract, the turns they have seen, the tools they have run, the summaries they
have written. How well an agent answers, how fast it responds, and how much it costs are all decided by
one question asked of every memory item: \textbf{in what form should it be kept, where should it live,
and when should it be (re)materialized?} Concretely, each item can be stored as plaintext, a vector, a
graph, a latent/KV activation, or baked into parameters; it can sit on GPU, CPU, disk, or a remote tier;
and it can be written inline while serving, consolidated during an idle window, or left untouched.

Today this single decision is \textbf{fragmented across three research communities}, each of which
optimizes one axis and freezes the others. Memory-operating-system work---MemOS~\cite{F1_li2025memos},
MemGPT~\cite{P14_packer2023memgpt}---provides the primitives to migrate an item between representations
and tiers, but \emph{decides when to migrate by hand-tuned heuristics}. Reinforcement-learned memory
managers---Memory-R1~\cite{F6_yan2025memoryr1}, Memory-as-Action~\cite{F7_zhang2025memoryasaction}---learn
\emph{what content} to add, update, or delete, but operate over a single fixed representation and never
choose a tier or schedule an offline write. Systems-level cache tiering---LMCache~\cite{F5_liu2025lmcache}%
---places bytes across GPU/CPU/disk/remote to minimize \emph{miss rate or time-to-first-token}, with no
notion of the downstream task utility that the memory actually serves.

Our central observation is that \textbf{these three axes are coupled and cannot be optimized separately.}
Promoting an item to a latent or parameter representation is an expensive offline write, but it is then
cheap and fast to serve; leaving it as plaintext on disk is cheap to write but slow to read and heavy on
context tokens. The best representation depends on the tier and on how often the item will be reused; the
best tier depends on the representation's serving cost and the item's size; and whether to pay the
promotion cost at all depends on the budget. This is precisely a \textbf{budget-constrained scheduling
problem}---and a large body of ML-for-systems work (Belady-oracle imitation for cache
replacement~\cite{S2_liu2020parrot}, GNN schedulers for clusters~\cite{S5_mao2019decima}, offline-RL cache
eviction in production~\cite{S4_coldrl2025}) demonstrates that such problems are learnable.

The reason to act now is that every substrate exists but nothing orchestrates them: MemCube migration, RL
memory managers, tiered KV caches, latent memory~\cite{F2_wang2025mplus}, and sleep-time
consolidation~\cite{F3_lin2025sleeptime} all landed by late 2025, unconnected. What is missing is a
\emph{learned scheduler} that, per item and under a shared budget, trades accuracy against latency and
cost, with \textbf{task utility---not miss rate---as the objective.}

We present \textbf{STOA} (Storage \& Tiered Orchestration for Agents), a learned control plane that
decides \textbf{representation $\times$ tier $\times$ timing jointly}, per memory item, as a single
budget-constrained Markov decision process. The operator sets a latency/cost/token budget $(L,C,B)$ at
serving time and one policy adapts, tracing the accuracy--latency--cost frontier without retraining. To
tame the combinatorics (fifty legal actions per item, across many items) STOA factors the policy into
representation, tier, and timing heads over a graph network of the memory store, and evaluates it only
over the legal joint actions so masking is exact. It warm-starts the policy by imitating a Belady oracle
and assigns credit for delayed, query-time payoffs with a mix of oracle imitation and process
supervision~\cite{F20_xiong2025raggym}.

Having built the formulation, we take seriously the empirical question it invites---\emph{is this decision
actually learnable?}---and answer it on production KV traces rather than on synthetic workloads. The
answer turned out to depend less on the model than on \emph{when the model is asked}.

Placement studies of this kind, including three revisions of this one, fix a split instant: observe a
prefix, build features from it, place every block, and bill the result against the future. Feature
construction can only describe a block that has been seen --- and on these traces 45--46\% of blocks have
no pre-split access at all, while the hindsight oracle they are scored against ranks them anyway. Those
blocks carry \textbf{97\% of the oracle's advantage}. A policy confined to the remaining 3\% looks flat
however well it ranks, which is precisely what our earlier ``ranking does not convert into placement''
result measured. We retract it.

Moving the decision to each block's \emph{arrival} --- placing it the moment it first appears, from the
session context available at that instant --- raises the reachable share from 3\% to \textbf{83.5\%} on
one trace and 2\% to \textbf{19.0\%} on the other. What captures it is not prediction:
first-come-first-served admission with a fixed action reaches 87.4\% and 50.2\% of the ceiling, while a
learned arrival-time ranker at AUC 0.574 and 0.635 reaches 84.5\% and 47.2\% --- less, on both. But the
arrival \emph{order} is not free: shuffling it costs 3.3 points on one trace and \textbf{14.1} on the
other. So the design implication is narrower than ``deciding at all is enough'': build an
\textbf{admission point} rather than a predictor, and do not assume the order it admits in is
incidental.

Reaction is competitive for a related reason. Tiering rules that forecast nothing at all --- LRU, LFU and
a fill-once cache --- reach 75--99\% of Belady's hit rate on one Mooncake trace and 40--94\% on the other,
and the whole adaptive-replacement family (ARC, LeCaR, CACHEUS) beats the best of them by \emph{at most
4.6 points}. All of them adapt on a ghost hit, and 76--78\% of blocks here are never reused, so most
evictions are correct and emit no signal: ARC gets one on 3.5--7.3\% of misses. A population of
singletons starves the channel adaptive replacement runs on.

We report how much measurement discipline this took, because it is the part most likely to transfer, and
because it is not a flattering story. Three claims were retracted, the last two only under adversarial
review \emph{after} the error catalogue below had been written: a headline that randomized sampling
showed had never been measurable, the decoupling result above, and an LRB column that turned out to
measure random eviction because a training buffer was truncated to its tail. Two further sub-claims of
the present framing were falsified later still, by a re-verification pass run after this paper was
rewritten: three of the artifacts it cited predated a fix to the code that produced them. In every case
the aggregate looked reasonable, which is the point --- none of our checks asked whether a reported
quantity was arithmetically able to be what it claimed, and none asked whether a stored result was
older than the program that wrote it.

\paragraph{Contributions.}
\begin{itemize}
  \item \textbf{(C1)} A unifying formulation of agent-memory placement as a budget-constrained MDP over
  representation $\times$ tier $\times$ timing, with a single serve-time $(L,C,B)$ knob
  (\Cref{sec:formulation}), and a control-plane architecture for it (\Cref{sec:arch}).
  \item \textbf{(C2)} A timing result on production traces. The split-instant protocol these studies use
  makes 97\% of a hindsight oracle's advantage unreachable, because 45--46\% of blocks do not exist yet
  when it decides (\Cref{sec:reachable}). Deciding at each block's \emph{arrival} raises the reachable
  share from 3\% to 83.5\% and 2\% to 19.0\%; there, first-come-first-served admission captures 87.4\%
  and 50.2\% of what is reachable and a learned ranker captures less than FCFS on both, while the arrival
  order is worth 14 points on one trace (\Cref{sec:arrival}). Once the metric is repaired, ranking quality
  converts monotonically on one trace and not the other, and never enough to help
  (\Cref{sec:capacity}).
  \item \textbf{(C3)} Measurements of reactive tiering on both sources against the adaptive-replacement
  family (ARC, LeCaR, CACHEUS) and a reimplemented LRB --- the best of them beating the best plain rule by
  at most 4.6 points of Belady across eight operating points --- together with an eviction-signal inversion
  that makes a policy tuned on one source a poor default for the other, and a measured explanation: KV
  traces are dominated by one-time blocks, which starve the ghost-hit channel all four adapt on
  (\Cref{sec:reaction}).
  \item \textbf{(C4)} A measured separation of placement from retrieval on LoCoMo at the sample size a
  power analysis demanded (300 questions per arm, paired tests): retrieval leads by a factor of four in
  10 of 10 dialogues, while placement still beats the same-regime embedding baseline
  (\Cref{sec:placement-vs-retrieval}).
  \item \textbf{(C5)} A first probe of the representation axis on real data --- three textual
  fidelities, not the vector, graph, latent or parameter forms of \Cref{sec:formulation}, which we never
  instantiate --- pointing \emph{against} our own framing: one representation leads at every budget
  rather than crossing. Reported as a direction with the sample size that would settle it, since none of
  the eight comparisons survives correction (\Cref{sec:representation}).
  \item \textbf{(C6)} A reusable evaluation protocol: sixteen measurement defects that each manufactured
  a plausible result here, grouped into six mechanisms, and the cheap controls that catch them ---
  including the rule that a null is reportable only alongside the sample size that would have resolved
  it, and the rule that every edge of the source $\rightarrow$ artifact $\rightarrow$ figure
  $\rightarrow$ paper chain needs its own freshness check (\Cref{sec:protocol}).
\end{itemize}

We position STOA against prior work along the three axes it unifies in \Cref{tab:positioning}; no prior
system occupies their intersection---deciding representation, tier, \emph{and} timing, learned, under a
shared budget, with task utility as the objective.

% Positioning matrix. \checkmark = learned/decided; $\circ$ = fixed/heuristic; -- = not applicable.
\begin{table}[t]
  \centering
  \small
  \setlength{\tabcolsep}{4pt}
  \begin{tabular}{@{}lcccccc@{}}
    \toprule
    System & Repr. & Tier & Timing & \makecell{Task\\util.} & \makecell{Budget\\cond.} & \makecell{Joint\\credit} \\
    \midrule
    MemOS~\cite{F1_li2025memos}            & $\circ$ & $\circ$ & $\circ$ & $\circ$ & $\circ$ & -- \\
    Memory-R1~\cite{F6_yan2025memoryr1}    & -- & -- & \checkmark & \checkmark & $\circ$ & $\circ$ \\
    Mem-as-Act.~\cite{F7_zhang2025memoryasaction} & -- & -- & \checkmark & \checkmark & $\circ$ & \checkmark$^\dagger$ \\
    LMCache~\cite{F5_liu2025lmcache}       & -- & \checkmark & $\circ$ & $\circ$ & $\circ$ & -- \\
    Cold-RL~\cite{S4_coldrl2025}           & -- & \checkmark & -- & $\circ$ & $\circ$ & \checkmark$^\dagger$ \\
    Sleep-time~\cite{F3_lin2025sleeptime}  & -- & -- & \checkmark & \checkmark & $\circ$ & -- \\
    \textbf{STOA (proposed)}               & $\ast$ & $\ast$ & $\ast$ & $\ast$ & $\ast$ & $\ast$ \\
    \bottomrule
  \end{tabular}
  \caption{No prior system decides representation, tier, \emph{and} timing jointly---learned, under a
  shared budget, with task utility as the objective. \checkmark: learned; $\circ$: fixed/heuristic; --:
  not applicable. $\dagger$: credit for one axis only (content or eviction). \textbf{$\ast$: proposed in
  \Cref{sec:formulation}--\ref{sec:arch} and \emph{not evaluated in this paper}} --- the row states an
  intended design, not a measured capability, and \Cref{sec:evidence} gives reasons to doubt parts of it.}
  \label{tab:positioning}
\end{table}

